\documentclass[11pt,a4paper]{article}
% ===================================================================
%  Vorlage: Karte aus Geodaten (Bundesländer, Nachbarn, Flüsse, Seen)
%  1. Daten erzeugen (siehe SKILL.md):  python3 geo2tikz.py ... -o karte_daten.tex
%  2. pdflatex vorlage_karte.tex
% ===================================================================
\usepackage[T1]{fontenc}
\usepackage{lmodern}
\usepackage[ngerman]{babel}
\usepackage[a4paper,margin=2cm]{geometry}
\usepackage{tikz}
\pagestyle{empty}
\setlength{\parindent}{0pt}

\input{karte_daten}

% ---- Aussehen: erst die ganze Ebene, dann einzelne Objekte ----------
\colorlet{wasser}{blue!45!cyan}
\tikzset{
  geo nachbarn/.style  = {fill=gray!12, draw=gray!50, line width=0.3pt},
  geo laender/.style   = {fill=green!8!white, draw=gray!70, line width=0.4pt},
  geo laender/DE-SN/.style = {fill=orange!45},          % Sachsen hervorheben
  geo fluesse/.style   = {draw=wasser, line width=0.6pt, line cap=round, line join=round},
  geo seen/.style      = {fill=wasser!40, draw=wasser, line width=0.3pt},
}

\begin{document}
{\Large\bfseries\sffamily Deutschland – Bundesländer und Flüsse}

\medskip
\begin{center}
\begin{tikzpicture}
  \clip (0,0) rectangle (12,15.2);        % Größe steht in karte_daten.tex
  \geonachbarn
  \geolaender
  \geoseen
  \geofluesse
  \node[font=\small\bfseries] at (laender-DE-SN) {SN};
  \node[font=\small] at (laender-DE-BY) {BY};
  \node[font=\footnotesize\itshape, text=wasser!80!black, fill=white, inner sep=1pt, opacity=.8, text opacity=1]
       at (fluesse-Elbe) {Elbe};
\end{tikzpicture}
\end{center}

\vfill
{\footnotesize Kartendaten: Natural Earth (gemeinfrei), naturalearthdata.com}
\end{document}
